\section{Singleton coordinates with zero advice slopes}
\label{sec:zero-advice}

Proposition~\ref{prop:matrix-decoding} converts high acceptance of the
matrix test into agreement between two local strategies supplied with
correlated linear advice.  We now bound the value of this advice
experiment on a NO parity instance.  The use of vanishing advice is part
of the smooth outer-game framework of
\cite[Sections~3.3--3.4]{KMS17}.  We give the full local reconstruction
argument for the present occurrence-indexed questions.  The key is to
identify singleton coordinates on which both advice slopes vanish.
After fixing the data outside those coordinates, their question pairs
have the ordinary independent equation--variable law.  Each side can
then reconstruct its original advice input from its own remaining
questions and the fixed data.  This defines strategies for an ordinary
parallel repetition, whose value bounds the original agreement.

We prove Lemma~\ref{lem:zero-advice} for the experiment in
Definition~\ref{def:projected-advice}, with $k>1$ an integer cube and
$\beta=k^{-2/3}$.  The distinct variable IDs in each equation make its
ordinary equation--variable game $G_{\mathcal E}$ a projection game;
Equation~\eqref{eq:base-game-gap} gives its value at most $14/15$ when
$\OPT(\mathcal E)\le4/5$.

\begin{proof}[Proof of Lemma~\ref{lem:zero-advice}]
Fix an arbitrary pair of local randomized strategies in the advice
experiment.  We first condition on $A$ and combine the two visible
advice maps into
\[
 D_O=(T_O,S_O):L_O\longrightarrow E_A,
 \qquad E_A=W\oplus\im A.
\]
Conditional on the questions and projection choices, the intercept
$c_O$ and all block slope columns of $D_O$ are independent uniform
elements of $E_A$.  Indeed, the columns of $T_O$ and $M_O$ are
independent and uniform, and applying $A$ to a uniform column gives a
uniform element of $\im A$.  This coefficient law depends only on
whether a block is an equation or a singleton, and not on any equation
ID, variable ID, or selected position.  In particular, $c_O$ is
independent of all block data.

Let $I\subseteq[k]$ consist of the coordinates $i$ projected to a singleton
whose advice slope $d_i\in E_A$ is zero.  Conditional on $A$,
the indicators of membership in $I$ are independent, with common probability
\begin{equation}
 p_A=\frac{\beta}{\abs{E_A}}
     =\beta\,2^{-\dim W-\rank A}.
 \label{eq:zero-advice-probability}
\end{equation}
More precisely, for every equation occurrence $e$ and position
$j\in[3]$,
\begin{equation}
 \Pr[e_i=e,\ i\text{ is a singleton at position }j,\ d_i=0\mid A]
 =\omega_e\,\frac{\beta}{3\abs{E_A}}.
 \label{eq:zero-advice-joint-law}
\end{equation}
Thus, conditional on belonging to $I$, a coordinate has precisely the
ordinary equation--variable distribution.

For the analysis, let $R$ consist
of $A$, $I$, the original intercept $c_O$, and the following data for
every $i\notin I$: its equation occurrence, whether the block was
retained or projected, the selected position in the latter case, and
its advice slope columns.  No witness, decoder output, agreement event,
or response randomness is included in $R$.

For each value $R=r$ of positive probability, regard the set
$I$ recorded in $r$ as fixed.  When $I$ is nonempty, we will use $r$
as constant data when defining two strategies for
$G_{\mathcal E}^{\otimes\abs{I}}$.
Each strategy will reconstruct one original advice input from its own
product-game questions, apply the corresponding original response rule,
and retain the answers in the slots $I$.

Conditional on $R=r$, the pairs
\[
 (e_i,v_{e_i,j_i}),\qquad i\in I,
\]
are independent question pairs of $G_{\mathcal E}$.  Indeed, conditional
on $A$, the coordinate data are independent, and
Equation~\eqref{eq:zero-advice-joint-law} shows that conditioning on $I$
preserves the ordinary law in each coordinate in $I$.  The outside
coordinate data and the independent intercept impose no further
condition on these pairs.

We now verify the two reconstructions, including the affine intercept
shared across the blocks.  An equation with right-hand side $b_e$ has actual
bits $x,y,b_e t_0+x+y$.  Pulling back a singleton slope $d_i$ thus
produces the two slopes
\begin{equation}
 (d_i,0),\quad (0,d_i),\quad (d_i,d_i)
 \label{eq:zero-advice-padded-slopes}
\end{equation}
at positions one, two, and three, respectively.  The intercept of
$D_U=D_O\bar\pi$ is
\begin{equation}
 c_U=c_O+
 \sum_{\substack{i\text{ projected to a singleton}\\j_i=3}}
 b_{e_i}d_i.
 \label{eq:zero-advice-intercept}
\end{equation}
For a coordinate in $I$, every slope in
Equation~\eqref{eq:zero-advice-padded-slopes} is zero and its contribution to
Equation~\eqref{eq:zero-advice-intercept} is zero.  All remaining terms in
that equation are determined by $r$.

The equation side receives its occurrences $e_i$ for $i\in I$.
Together with $r$, they specify all of $U$.  Its two advice slopes in
each block in $I$ are zero, while $r$ determines every other slope and
the full intercept.  These coefficients reconstruct both padded advice
maps without using any selected position or variable question in $I$.
The variable side receives the variable question for each $i\in I$.
Together with $r$, these specify all of $O$.  Its singleton
slopes in $I$ are zero, and $r$ supplies every other slope and the original
intercept.  These coefficients reconstruct its two advice maps without
using any equation in $I$.  Repeated IDs across different blocks cause no
additional restriction: questions were sampled independently with
replacement, and answers use the local copies specified in
Section~\ref{sec:game}.

These reconstructions define the promised product-game strategies:
each side applies its original response rule to the reconstructed
input and retains its answers in the slots $I$, in their original
order.  Their equation answers are satisfying triples because the
original answers are valid.  For every realization of the questions in $I$ and response randomness, full agreement in the advice
experiment implies acceptance in every slot in $I$.  If $I=\varnothing$,
the conditional success probability is at most one.  Otherwise,
\[
 \Pr[\bar\pi a_U=a_O\mid R=r]
 \le \val(G_{\mathcal E}^{\otimes\abs{I}})
 \le \left(1-\frac{1}{3600}\right)^{\abs{I}}.
\]
The first inequality uses the strategies just constructed and the
conditional product law.  The second uses
Theorem~\ref{thm:repetition} and Equation~\eqref{eq:base-game-gap}.
The original response rules are evaluated only at their original
local inputs; $r$ indexes the auxiliary product-game strategies.
Thus the final bound holds for every $r$, including $I=\varnothing$
when its right-hand side is interpreted as one.

Finally, conditional on $A$, the size of $I$ has distribution
$\operatorname{Binomial}(k,p_A)$.  Averaging the preceding inequality
first over $R$ conditional on $A$, and then over $A$, gives
\begin{align*}
 \Pr\bigl[\bar\pi(a_U)=a_O\bigr]
 &\le \E_A\left(1-\frac{p_A}{3600}\right)^k\\
 &\le\left(1-\frac{\beta\,2^{-\dim W-r_{\rm s}}}{3600}\right)^k\\
 &\le \exp\!\left(
 -\frac{k\beta\,2^{-\dim W-r_{\rm s}}}{3600}
 \right),
\end{align*}
since $\rank A\le r_{\rm s}$.  Substituting
$k\beta=k^{1/3}$ proves Equation~\eqref{eq:zero-advice-agreement-bound}.
\end{proof}
